\section{Step 15 Form capital sigma zero of old a}\label{sha-s15}

Rotate old a by two, thirteen and twenty-two bits and XOR those words into \texttt{t1}. Retain T1 in \texttt{t0}. Two distinct destination temporaries prevent the second calculation from destroying the first. The next region will add the majority term to \texttt{t1}.

\subsection{Write and inspect}
Replace only the temporary instructions between \texttt{\# BEGIN comp\_sigma0} and \texttt{\# END comp\_sigma0}. Keep the supplied labels outside that region.

\Needspace{9\baselineskip}
\MARSCodeLabel{MIPS assembly}
\begin{minted}[samepage=true]{mars-idea}
    rotr $t1, $s0, 2
    rotr $t2, $s0, 13
    xor $t1, $t1, $t2
    rotr $t2, $s0, 22
    xor $t1, $t1, $t2
\end{minted}

Use recipe \texttt{s15-1} with the cumulative program through this edit. The exact recipe names the entry, prepared input and checkpoint. Observe the checkpoint before continuing to its temporary or normal return.

\subsection{Expected observation}
For the first abc round, \texttt{t1} has low word \texttt{ce20b47e} while \texttt{t0} still holds T1.

\Needspace{9\baselineskip}
\subsection{Check your understanding}
Which temporary must survive this calculation?

\begin{itemize}
\item \textbf{A} The already computed T1.
\item \textbf{B} The original source byte in \texttt{a0}.
\item \textbf{C} The current syscall number.
\end{itemize}
